% ----------------------------------------------------------
% Conclusão
% ----------------------------------------------------------
\chapter{Conclusão}
O desenvolvimento deste trabalho permitiu compreender, de forma integrada, tanto os aspectos jurídicos da execução penal quanto os fundamentos técnicos necessários para a modelagem de um sistema capaz de apoiar sua gestão de maneira eficaz. A análise do fluxo da execução penal demonstrou que a administração correta de informações — como remição, detração, unificação de penas, datas‑base e elegibilidade para benefícios — é essencial para assegurar a observância dos direitos da pessoa condenada e a atuação precisa dos operadores do Direito. Esse entendimento fundamentou a definição das funcionalidades do sistema e orientou a construção dos modelos conceituais apresentados.

A partir desse estudo, foi possível identificar as principais necessidades do processo de execução penal e traduzi‑las em funcionalidades objetivas: registrar remição, registrar detração, registrar unificação de pena, calcular benefícios, exibir status de progressão, processar progressão de regime, além do cadastro de condenados e advogados. Essas funcionalidades estruturam o núcleo operacional do sistema e refletem diretamente a dinâmica prevista na Lei de Execução Penal (LEP), permitindo o acompanhamento claro, auditável e automatizado do cumprimento da pena.

A modelagem proposta — por meio dos diagramas de casos de uso, diagrama de classes e do diagrama Entidade‑Relacionamento (DER) no estilo Chen — consolidou a visão lógica e estrutural do sistema. Os diagramas demonstram como cada entidade se relaciona com o processo jurídico real, garantindo coerência e fidelidade ao funcionamento da execução penal. A criação do DER, em especial, mostrou como a informação é organizada, preservada e cruzada para suportar o cálculo dos benefícios e a geração do status de progressão, elementos centrais para a atuação da Vara de Execuções Penais.

Desse modo, o sistema concebido neste trabalho não é apenas uma ferramenta tecnológica, mas uma solução alinhada às demandas práticas do cotidiano forense, contribuindo para maior segurança jurídica, eficiência administrativa e transparência no controle da pena. Além de favorecer decisões mais rápidas e precisas, a modelagem apresentada coloca o condenado no centro do processo de execução, garantindo que seus direitos sejam acompanhados de forma objetiva, sistemática e verificável.

Portanto, conclui-se que o sistema proposto atende plenamente aos requisitos identificados, oferecendo uma base sólida tanto para implementação prática quanto para evolução futura. A integração entre teoria jurídica e modelagem computacional demonstrou que é possível transformar estruturas complexas da execução penal em fluxos claros e automatizáveis, contribuindo para uma justiça mais acessível, eficiente e humanizada.
